%
% Der Satz von Rothstein-Trager
%
\subsection{Der Satz von Rothstein-Trager}
Die Form des logarithmischen Teils des Integrals ist zwar bekannt,
und der Satz~\ref{buch:ratint:rothsteintrager:satz:residuum}
ermöglicht eine direkte Berechnung des Koeffizienten, sobald
die Nullstellen bekannt sind, aber genau die Bestimmung der
Nullstellen war ja der Punkt, warum die Methode von Bernoulli
als unzureichend erkannt wurde.
Dies ändert der folgende Satz von Rothstein-Trager.

\begin{satz}[Rothstein-Trager]
\label{buch:ratint:rothsteintrager:satz:rothsteintrager}
\index{Rothstein, Michael}%
\index{Trager, Barry M.}%
Ist $p(z)$ teilerfremd zum normierten, quadratfreien Polynom $q(z)$ mit
$\deg p(z)<\deg q(z)$, dann ist
\begin{equation}
\int
\frac{p(z)}{q(z)}
\,dz
=
\sum_{k=1}^m c_k \log v_k(z),
\label{buch:ratint:rothsteintrager:eqn:rtintegral}
\end{equation}
wobei die $c_k$ die verschiedenen, möglicherweise komplexen Zahlen $c$ sind
derart, dass $p(z)-cq'(z)$ und $q(z)$ einen nichttrivialen gemeinsamen
Teiler haben.
Die zugehörigen $v_k(z)$ sind
\[
v_k(z)
=
\operatorname{ggT}\bigl(p(z)-c_kq'(z),q(z)\bigr).
\]
Sie sind normiert, teilerfremd und quadratfrei.
\end{satz}

Aus der Methode von Bernoulli ist bekannt, dass die Stammfunktion tatsächlich
\index{Bernoulli!Methode von}%
die Form~\eqref{buch:ratint:rothsteintrager:eqn:rtintegral} hat.
Danach sind die $v_k(z)$ verschiedene Linearfaktoren, deren Produkt
\[
q(z)
=
\prod_{k=1}^m v_k(z)
\]
ist.
Falls zwei Faktoren $v_i(z)$ und $v_k(z)$ den gleichen Koeffizienten
$c=c_i=c_k$ haben, können die beiden Terme zu einem Term $v(z)$
zusammengefasst werden mit $v(z) = v_i(z)v_k(z)$.
Es ist nämlich
\[
cv(z)
=
c\log \bigl(v_i(z)v_k(z)\bigr)
=
c(\log v_i(z) + \log v_k(z))
=
c_i \log v_i(z) + c_k\log v_k(z).
\]
Man darf also sogar davon ausgehen, dass die $c_k$ alle verschieden
sind und die $v_k(z)$ Produkte von verschiedenen Linearfaktoren.
Somit haben die $v_k(z)$ nur einfache Nullstellen, sie sind also
alle quadratfrei.

\begin{proof}
Nach der Vorbemerkung muss noch gezeigt werden, dass die Konstanten
$c_k$ und die Polynome $v_k(z)$ auf die genannte Art gefunden werden
können.
Wir nehmen dazu an, dass die Stammfunktion die Form 
\eqref{buch:ratint:rothsteintrager:eqn:rtintegral} hat
und schreiben
\begin{align*}
u_k(z) &= \prod_{\scriptstyle i=1\atop \scriptstyle i\ne k}^m v_i(z)
&&\text{und}&
v(z)   &= \prod_{i=1}^m v_i(z).
\end{align*}
Die Ableitung von \eqref{buch:ratint:rothsteintrager:eqn:rtintegral} ist
\[
\frac{p(z)}{q(z)}
=
\sum_{k=1}^m c_k \frac{v'_k(z)}{v_k(z)}.
\]
Durch Multiplizieren mit $q(z)$ und allen $v_i(z)$, $i=1,\dots,m$,
erhalten wir
\begin{equation}
p(z)\prod_{i=1}^m v_i(z)
=
\sum_{k=1}^m
c_k
q(z)
v'_k(z)
\prod_{\scriptstyle i=1 \atop \scriptstyle i\ne k}^m v_i(z)
=
\sum_{k=1}^m
c_k
q(z)
v'_k(z)
u_k(z).
\label{buch:ratint:rothsteintrager:eqn:produkt}
\end{equation}

Für den Beweis gehen wir in den folgenden Schritten vor:
\begin{enumerate}
\item Schritt:
$q(z)$ teilt das Produkt $v(z)$ aller $v_i(z)$.
Die rechte Seite von \eqref{buch:ratint:rothsteintrager:eqn:produkt}
enthält den Faktor $q(z)$, also muss auch die linke Seite durch $q(z)$
teilbar sein.
Da $p(z)$ und $q(z)$ teilerfremd sind, muss $q(z)$ ein Teiler des
Produktes auf der linken Seite von
\eqref{buch:ratint:rothsteintrager:eqn:produkt} sein.

\item Schritt:
Jeder Faktor $v_i(z)$ teilt $q(z)$.
$v_i(z)$ ist ein Faktor der linken Seite von
\eqref{buch:ratint:rothsteintrager:eqn:produkt}
und jedes Faktors $c_kq(z)v'_k(z)u_k(z)$ der rechten Seite ausser
im Fall $k=i$.
Daher muss $v_i(z)$ ein Teiler von $c_iq(z)v'_i(z)u_i(z)$ sein.
Die Polynome $v_i(z)$ und $u_i(z)$ sind teilerfremd und weil $v_i(z)$
quadratfrei sind auch $v_i(z)$ und $v'_i(z)$.
Somit muss $v_i(z)$ ein Teiler von $q(z)$ sein.

\item Schritt:
Jeder Faktor $v_i(z)$ ein Teiler von $q(z)$, also ist auch
das Produkt $v(z)$ ein Teiler von $q(z)$.
Da sowohl $v(z)$ als auch $q(z)$ normierte Polynome sind, folgt aus
$v(z)\mid q(z)$ und $q(z)\mid v(z)$, dass $v(z)=q(z)$.

\item Schritt: wegen $v(z)=q(z)$ kann man jetzt beide Seiten von
\eqref{buch:ratint:rothsteintrager:eqn:produkt}
durch $q(z)$ teilen und erhält
\begin{equation}
p(z)
=
\sum_{i=1}^m
c_i
v'_i(z)
u_i(z).
\label{buch:ratint:rothsteintrager:eqn:produkt2}
\end{equation}

\item Schritt:
Die Gleichung
\eqref{buch:ratint:rothsteintrager:eqn:produkt2} kann auch als
\[
p(z)-c_k v'_k(z) u_k(z)
=
\sum_{\scriptstyle i=1\atop \scriptstyle i\ne k}^m c_iv'_i(z) u_i(z)
\]
geschrieben werden.
Ein nichttrivialer Teiler $t(z)$ von $v_k(z)$ teilt alle
Summanden auf der rechten Seite und muss daher auch ein Teiler
von $p(z)-c_k v'_k(z) u_k(z)$ sein.
Insbesondere ist $v_k(z)$ selbst ein Teiler von $p(z)-c_kv'_k(z)u_k(z)$
oder anders ausgedrückt: $v_k(z)$ ist der grösste gemeinsame Teiler von
$p(z)-c_kv'_k(z)u_k(z)$ und $v_k(z)$.

\item Schritt:
Ein nichttrivialer Teiler $t(z)$ von $v_k(z)$ teilt jeden Term $u_i(z)$
mit $i\ne k$, somit ist $t(z)$ ein Teiler von
\[
p(z) - c_kv'_k(z) u_k(z)
-
c_k \sum_{\scriptstyle i=1\atop\scriptstyle i\ne k}
v'_i(z)u_i(z)
=
p(z) - c_k \sum_{i=1}^m v'_i(z)u_i(z)
=
p(z) - c_k q'(z)
\]
sein.

\item Schritt:
Ein nichttrivialer Teiler von $v_i(z)$ mit $i\ne k$ ist ein Teiler von
$p(z)-c_iq'(z)$.
Wäre er auch ein Teiler von $p(z)-c_kq'(z)$, dann müsste er auch ein Teiler
der Differenz
\[
p(z)-c_iq'(z)-(p(z)-c_kq'(z))
=
(c_i-c_k)q'(z)
\ne 0
\]
sein.
Dann wäre er aber auch ein gemeinsamer Teiler von $q(z)$ und $q'(z)$, was
nicht möglich ist, weil $q(z)$ quadratfrei ist.

\item Schritt:
Ein gemeinsamer Teiler von $p(z)-c_kq'(z)$ und $b(z)$ kann in Faktoren
$t_i(z)$ zerlegt werden, die $p(z)-c_kq'(z)$ und $v_i(z)$ teilen.
Nach dem letzten Schritt müssen die Faktoren $t_i(z)$ mit $i\ne k$
trivial sein.
Ist ein gemeinsamer Teiler von $p(z)-c_kq'(z)$ und $b(z)$ auch
ein gemeinsamer Teiler von $p(z)-c_kq'(z)$ und $v_k(z)$.
Insbesondere ist $v_k(z)$ der grösste gemeinsame Teiler von $p(z)-c_kq'(z)$
und $q(z)$.
\end{enumerate}
Damit ist alles bewiesen.
\end{proof}

Wenn es gelingt, eine algebraische Methode zur Bestimmung der
Koeffizienten $c_i$ zu finden, können auch die Faktoren 
$r_i(z)$ gefunden werden.
Man beachte, dass nicht mehr eine Faktorisierung von $q(z)$
verlangt wird, sondern nur noch die Möglichkeit, gemeinsame
Faktoren zu finden, was im Prinzip mit dem euklidischen Algorithmus
möglich ist.

\begin{beispiel}
Der Integrand von
\[
\int \frac{1}{z^3+z}
\,dz
\]
hat nicht die Form, die man in der Methode von Bernoulli erwarten
würde, aber der Nenner ist quadratfrei, die Methode von Rothstein-Trager
ist daher anwendbar mit $p(z)=1$ und $q(z) = z^3+z$.
Es müssen daher $c$ gefunden werden, so dass 
\[
1-c\cdot (3z^2+1)
\qquad\text{und}\qquad
z^3+1
\]
einen nichttrivialen gemeinsamen Teiler haben.
Der euklidische Algorithmus ergibt
\begin{align}
z^3+z
&=
-\frac{1}{3c}z
\cdot
(1-c - 3cz^2)
+
\frac{(2c+1)z}{3c}
\label{buch:ratint:rothsteintrager:bsp1:rest1}
\\
(1-c - 3cz^2)
&=
-\frac{9c^2z}{2c+1}
\cdot
\frac{(2c+1)z}{3c}
+
1-c
\label{buch:ratint:rothsteintrager:bsp1:rest2}
\end{align}
In \eqref{buch:ratint:rothsteintrager:bsp1:rest1} verschwindet
der Rest, falls $c_1=-\frac12$ ist.
Das zugehörige $v_1(z)$ erhält man, indem man $c_1$ in den Divisor
\[
1-c_1-3c_1z^2
=
1+\frac12+\frac32z^2
=
\frac32(z^2+1)
\qquad\Rightarrow\qquad
v_1(z)
=
z^2+1
\]
einsetzt.
Nimmt man $c\ne -\frac12$ an, folgt aus dem zweiten
Schritt
\eqref{buch:ratint:rothsteintrager:bsp1:rest2}
des euklidischen Algorithmus, dass ein nichttrivialer Teiler vorliegt,
wenn $1-c=0$ ist, also $c_2=1$.
Einsetzen von $c_2$ in den Divisor ergibt
\[
\frac{(2c_2+1)z}{3c_2}
=
\frac{3z}{3}
=
z
=
v_2(z).
\]
Damit kann man jetzt das Integral
\[
\int \frac{1}{z^3+z}\,dz
=
c_1\log v_1(z)
+
c_2\log v_2(z)
=
-\frac12\log(z^2+1)
+
\log z
\]
hinschreiben.
\end{beispiel}

\begin{beispiel}
Man wende die Methode von Rothstein-Trager auf das Integral
\[
\int \frac{1}{z^2+1}\,dz
\]
an.

\smallskip

\noindent
In diesem Beispiel ist $p(z) = 1$ und $q(z)=z^2+1$.
Es müssen die Konstanten $c_k$ gefunden werden, für die
\[
p(z)-c\cdot q'(z)
=
1-c\cdot 2z
\qquad\text{und}\qquad
q(z) = z^2 +1
\]
einen nichttrivialen gemeinsamen Teiler haben.
Aus der Faktorisierung $z^2+1=(z+i)(z-i)$ kann man ablesen, dass
\[
1-c\cdot 2z
=
-2c\biggl(z-\frac{1}{2c}\biggr)
=
-2c(z\pm i)
\]
sein muss, woraus man
\[
-\frac{1}{2c}=\pm i
\qquad\Rightarrow\qquad
c
=
\pm\frac12i
\]
ablesen kann.
Es folgt, dass
\[
\int \frac{1}{z^2+1}\,dz
=
\frac{i}{2}\log (z+i)
-
\frac{i}{2}\log (z-i)
=
\frac{i}2
\log
\frac{z+i}{z-i}
=
\arctan z
\]
mit der aus der eulerschen Formel ableitbaren Darstellung der
Arkustangensfunktion durch komplexe Logarithmusfunktionen
(siehe dazu auch später Satz~\ref{buch:intalg:elementar:satz:trigo:tancot}).
\end{beispiel}

%
% Gemeinsame Nullstellen und die Resultante
%
\subsubsection{Gemeinsame Nullstellen und die Resultante}
Der Satz~\ref{buch:ratint:rothsteintrager:satz:rothsteintrager}
verlagert das Problem, die Nullstellen des Nenners zu finden darauf,
die Koeffizienten $c_i$ zu finden.
Dies ist aus zwei Gründen einfacher:
\begin{enumerate}
\item
Oft sind einige der Koeffizienten $A_k$ gleich.
In diesen Fällen ist es nicht nötig, die zugehörigen Linearfaktoren
zu kennen, es genügt das Produkt dieser Linearfaktoren zu finden,
dies ist das Polynom $r_i(z)$ im
Satz~\ref{buch:ratint:rothsteintrager:satz:rothsteintrager}.
\item
Für das Kriterium, einen gemeinsamen Teiler zu finden gibt es
ein einfaches algebraisches Kriterium, nämlich die Resultante.
\end{enumerate}

In Kapitel~\ref{chapter:resultante} wird beschrieben, wie die 
Resultante des Problem, die Koeffizienten $c_i$ zu finden, lösen
kann.
\index{Resultante}%

